\chapter{The geometric inequality II: local sizes, covolume and the Chebyshev inequality}\label{ch:geob}

This chapter covers the size data that the archive attaches to a package: the
archimedean main term (Stirling's $y\log y-y$ and a period $\log\Omega$ built
from $L'/L(1,\chi_{-d})$), the local sizes at finite places with the
Legendre-type count at inert primes, the two ``link'' predicates that bundle
these with the period lattices, and the covolume statements on the indefinite
side. It ends with the passage from a slope record to the gap bound
(\texttt{ch\_top}) and from tied families to the Chebyshev inequality on
windows (\texttt{level\_gm\_le\_on\_family\_scales}), together with the
statement \texttt{TopG} that the existence part of the archive must supply.
The lattices, slopes and the slope record itself are treated in the preceding
chapter and are referred to here only through their Lean names.

Throughout, $D$ is a package (the archive's \texttt{D7}) and $L$ a size datum
for $D$ (\texttt{SizeData₂ D}). The scale of $L$ has real numbers $x>1$ and
$0<\eps<1/52$, and the archive sets
\[
L_x=x^{\eps},\quad b=e^{x^{\eps}},\quad a=xb,\quad X'=e^{x^{\eps/3}},\quad
b'=a/X',\quad w''(n)=b'/(xn),\quad X''(n)=a/w''(n)=xnX'
\]
(\arch{Formalism/LocalRingCardinality/DataParams₂.lean}, lines 9--47).
We write $n$ for the order $|A|$ of the host's group, $f=2n$, $m$ for the
number of lifts and $k$ for the base rank of $D$.

\section{The archimedean side}

\begin{definition}[Stirling's function and the prime-sum comparison]\label{geob:phF}
\lean{Formalism.BasicNotions.phF, Formalism.BasicNotions.Gfun}
For real $y$ put $\operatorname{ph}(y)=y\log y-y$. Put
\[
G(y)=\sum_{1\le n\le\lfloor y\rfloor}\Lambda(n)\Bigl(\frac yn-1\Bigr)-y\log y+(1+\gamma)y,
\]
with $\Lambda$ Mathlib's von Mangoldt function and $\gamma$ Euler's constant.
\end{definition}

\begin{definition}[the quadratic character of the host and its logarithmic derivative]\label{geob:dlogL}
\lean{Formalism.BasicNotions.chiK, Formalism.BasicNotions.dlogL}
% uses-lean: Formalism.BasicNotions.D7, Formalism.BasicNotions.Host
\uses{host:host-full} % resolved
Let $d$ be the prime of the host of $D$ (so $K=\Q(\sqrt{-d})$ sits in the host
field $E$). Let $\chi_K$ be the Dirichlet character modulo $d$ obtained from
the quadratic character of $\Z/d\Z$ (the Legendre symbol modulo $d$), viewed
with values in $\CC$. For $s\in\CC$ put
\[
\operatorname{dlogL}(D,s)=\Re\frac{L'(s,\chi_K)}{L(s,\chi_K)},
\]
where $L(s,\chi_K)$ is Mathlib's $L$-function of $\chi_K$ and $L'$ its complex
derivative.
\end{definition}

\begin{definition}[root discriminant, $v$ and the period $\log\Omega$]\label{geob:logOm}
\lean{Formalism.BasicNotions.Host.rdE, Formalism.BasicNotions.Host.v, Formalism.BasicNotions.SizeData₂.logOm}
% uses-lean: Formalism.BasicNotions.SizeData₂
\uses{geob:dlogL}
For the host field $E$ put $\operatorname{rd}_E=|\operatorname{disc}E|^{1/[E:\Q]}$
and $v=\operatorname{rd}_E/2$. A size datum $L$ records a real number
$\ell_K$ (\texttt{logDerivativeAtOne}) together with the field
$\ell_K=\operatorname{dlogL}(D,1)$. Put
\[
\log\Omega=\gamma+\log\frac{\pi}{v}-\ell_K
=\gamma+\log\frac{2\pi}{\operatorname{rd}_E}-\Re\frac{L'(1,\chi_K)}{L(1,\chi_K)}.
\]
The symbol $\log\Omega$ is notation; the archive does not define $\Omega$ itself.
\end{definition}

\begin{definition}[the archimedean main term of a word]\label{geob:main}
\lean{Formalism.BasicNotions.SizeData₂.MAIN}
\uses{geob:phF, geob:logOm}
% uses-lean: Formalism.BasicNotions.SizeData₂
Let $TT$ be the set of top places of the host and $a_t$ ($t\in TT$) the integer
weights recorded by $L$; the size datum requires $a\le a_t\le a+a/(kx^{13})$.
For a word $u$ of $D$ let $T(u)$ be its level prime. Put
\[
\mathcal M(u)=\sum_{t\in TT}\Bigl(m\bigl(\operatorname{ph}(a_t)-a_t\log\Omega\bigr)-a_t\log T(u)\Bigr).
\]
\end{definition}

\begin{remark}[three functions called \texttt{MAIN}]
The archive has at least three unrelated declarations named \texttt{MAIN}: the
one above (\arch{Formalism/LocalRingCardinality/DataParams₂.lean}, line 232), a
complex-valued main term for one weight in
\arch{Formalism/SizeWindowBrackets/FirstPrimeNeSecondPrime.lean} (line 89),
and a third in \arch{Development/AutomorphicArithmetic/DedekindLFunctions/RunBaseDimensionLeRankAddGenusTerm.lean}
(line 110). There are also at least four declarations named \texttt{logOm},
all of the shape $\gamma+\log(\pi/v)-\ell$ for different records. Only the first
\texttt{MAIN} enters the upper link below.
\end{remark}

\begin{remark}[reading the main term]
Stirling's formula gives $\log\Gamma(y+1)=y\log y-y+O(\log y)$, so
$\operatorname{ph}(a_t)$ is the size of $\log a_t!$. The term
$-a_t\log\Omega$ normalises by the $a_t$-th power of a period whose logarithm
is $\gamma+\log(\pi/v)-L'/L(1,\chi_K)$, a formula of Chowla--Selberg or
Kronecker-limit type for CM by $K$, and the factor $m$ counts the lifts. This is
the only place in the size data where $L'(1,\chi_K)/L(1,\chi_K)$ enters. The
reading as a CM period is an interpretation; the archive only defines the
number.
\end{remark}

\section{Local sizes at finite places}

\begin{definition}[local sizes]\label{geob:sz}
\lean{Formalism.BasicNotions.lw, Formalism.BasicNotions.SizeData₂.lwP, Formalism.BasicNotions.SizeData₂.sz, Formalism.BasicNotions.SizeData₂.lifts, Formalism.BasicNotions.SizeData₂.gainD, Formalism.BasicNotions.SizeData₂.Fo, Formalism.BasicNotions.SizeData₂.Ga, Formalism.BasicNotions.SizeData₂.Ee}
% uses-lean: Formalism.BasicNotions.SizeData₂
Let $\mathbb L$ be the coefficient number field of $D$. For a finite place $w$
of $\mathbb L$ put $\operatorname{lw}_w=\log N(w)/[\mathbb L:\Q]$, let $e_w$ be
its ramification index over $\Z$, and for real $y$ put
$\operatorname{sz}_w(y)=2\operatorname{lw}_w\,e_w\,y$. The lift indices are
$j\in\{1,\dots,m\}$. A size datum carries a finite support set $S$ of places
and real functions $\operatorname{gain}_j(w)$, $F_j(w)$, $E(u,j,w)$,
$\operatorname{ta}_w$, $\operatorname{io}_w$ vanishing off $S$, and a set
$\mathrm{drop}$ of primes; $\operatorname{gainD}_j(w)$ is $0$ if the residue
characteristic of $w$ is in $\mathrm{drop}$ and $\operatorname{gain}_j(w)$
otherwise. Put
\[
\mathrm{Fo}=\sum_j\sum_{w\in S}\operatorname{sz}_w(F_j(w)),\quad
\mathrm{Ga}=\sum_{w\in S}\operatorname{sz}_w\Bigl(\sum_j\operatorname{gainD}_j(w)-\operatorname{ta}_w\Bigr),\quad
\operatorname{Ee}(u)=\sum_j\sum_{w\in S}\operatorname{sz}_w(E(u,j,w)).
\]
The size datum also forces $\operatorname{gain}_j(w)=0$ at the two auxiliary
primes $g_0,g_1$, and $\operatorname{gain}_j(w)=\operatorname{ta}_w=\operatorname{io}_w=0$
unless the residue characteristic of $w$ is inert in $K$ in the sense of
Definition~\ref{geob:inert}.
\end{definition}

\begin{definition}[prime classes and inert primes]\label{geob:inert}
\lean{Formalism.BasicNotions.PrimeClass, Formalism.BasicNotions.primeClass, Formalism.BasicNotions.Host.cls, Formalism.MainNotions.InertK}
% uses-lean: Formalism.BasicNotions.Host
\uses{host:host-full} % resolved
A host assigns to each natural number $p$ one of six classes: $\mathrm{ram2}$
if $p$ lies in a fixed finite set of small primes; $\mathrm{ramQ}$ if $p=q$;
$\mathrm{ramD}$ if $p=d$; otherwise $\mathrm{splitK}$ if the Jacobi symbol
$\bigl(\frac{-d}{p}\bigr)$ equals $1$; otherwise $\mathrm{signal}$ if the first
sign component of the Artin symbol of $p$ in the host is $+1$; and
$\mathrm{inertClass}$ in the remaining case. The prime $p$ is \emph{inert}
(\texttt{InertK}) if its class is $\mathrm{signal}$ or $\mathrm{inertClass}$.
\end{definition}

\begin{remark}[five predicates called \texttt{InertK}]
The name \texttt{InertK} has at least five declarations with different
bodies. The one in \arch{Formalism/SharpRestriction/GlobalTheoremInputs.lean}
(line 97) is the one above; the copy in
\arch{Formalism/SharpRestriction/GlobalSchurBounds/ZetaZeroFreeRegion.lean}
adds that $p$ is prime, and the one in
\arch{Formalism/SharpRestriction/GlobalSchurBounds/GL3Theta.lean} is
$\bigl(\frac{-d}{p}\bigr)=-1$. The lower link uses the first.
\end{remark}

\begin{definition}[the Legendre-type class sum]\label{geob:vm}
\lean{Formalism.BasicNotions.SizeData₂.vonMangoldtClassSum}
% uses-lean: Formalism.BasicNotions.SizeData₂
A size datum carries real numbers $w_{\mathrm{part}}(j,t,p,i)\ge1$. For a lift
index $j$, a top place $t$ and a natural number $p$ put
\[
\mathrm{vM}(j,t,p)=\sum_{\substack{0\le i<\lfloor\log_p a_t\rfloor+2\\ i\text{ odd}}}
\max\Bigl(0,\ \frac{a_t}{p^i}-w_{\mathrm{part}}(j,t,p,i)\Bigr),
\]
with $\lfloor\log_p a_t\rfloor$ Mathlib's \texttt{Nat.log}.
\end{definition}

\begin{remark}[Legendre's formula restricted to odd powers]
Legendre's formula is $v_p(N!)=\sum_{i\ge1}\lfloor N/p^i\rfloor$. The sum
$\mathrm{vM}$ keeps only odd $i$ and replaces $\lfloor N/p^i\rfloor$ by
$\max(0,a_t/p^i-w_{\mathrm{part}})$ with $w_{\mathrm{part}}\ge1$, which is at
most $\lfloor a_t/p^i\rfloor$ when $w_{\mathrm{part}}\ge1$. So
$\mathrm{vM}(j,t,p)\log p$ is a lower bound for the part of $\log_p$-valuation
of $a_t!$ coming from odd powers of $p$. Summed over inert $p$ the leading terms
give $\sum_{p\text{ inert}}(a_t/p)\log p$, a prime sum over a Frobenius class,
which against $\operatorname{ph}(a_t)$ is the archimedean-versus-finite balance
of $\log N!$ that the function $G$ of Definition~\ref{geob:phF} measures. This
reading is an interpretation of the definitions; the archive states the gain
bounds only as hypotheses (Definition~\ref{geob:lowerlink}).
\end{remark}

\begin{definition}[the lower-order terms of the lower link]\label{geob:loG}
\lean{Formalism.MainNotions.flG, Formalism.MainNotions.rG, Formalism.MainNotions.leastPowerExponent, Formalism.MainNotions.loG}
\uses{geob:inert}
For real $C$ and a natural number $p$ put
$\mathrm{fl}(p)=Cf\,N_m^2\log_p a+Cf(\log_p a)^2$, with $N_m=k+m(t+1)$ the
top rank of $D$. Put $r_G(\eps)=\lceil 8(4+8\eps)/\eps\rceil+1$ and
$e_p=\lceil\log_p(k^2x^{78})\rceil+1$. For real $C,\beta$ put
$\mathrm{lo}(p)=C(\mathrm{Sc}+k+\eta)\log_p(2a)+\mathrm{pr}(p)+\mathrm{or}(p)$,
where $\mathrm{Sc},\eta$ are fields of $L$, and $\mathrm{pr}(p)$,
$\mathrm{or}(p)$ vanish unless $p$ has class $\mathrm{inertClass}$ and
$p<X'$; in that case, with $p_0=k^2x^{26}$,
\[
\mathrm{pr}(p)=\begin{cases}Cr_G(\eps)\,b'\log x,&p<p_0,\\
C\bigl(b'/x+\eta+\mathrm{Sc}+(w''(n)+t)\log x/\delta_1\bigr)+\beta,&p_0\le p<X',\end{cases}
\]
with $\delta_1=x^{-\eps/8}$, and
$\mathrm{or}(p)=C\log x\bigl(\mathrm{La}_G\,a/(pk)+a\,p^{-e_p}\bigr)$ if
$p<p_0$, $\mathrm{or}(p)=C\log x\bigl(\mathrm{La}_G\,a/(pk)+a/p^2\bigr)$ if
$p_0\le p<X'$.
\end{definition}

\begin{definition}[the lower link]\label{geob:lowerlink}
\lean{Formalism.MainNotions.LinkG_at, Formalism.MainNotions.LinkG}
\uses{geob:sz, geob:inert, geob:vm, geob:loG}
% uses-lean: Formalism.MainNotions.G.Inst, Formalism.MainNotions.G.Consts, Formalism.MainNotions.G.VV, Formalism.MainNotions.G.Sv, Formalism.MainNotions.G.Nv, Formalism.MainNotions.G.ServedOrbitsContained, Formalism.MainNotions.G.MeanKappaBound, Formalism.MainNotions.FamP
\uses{geoa:famcore} % resolved
Let $X$ be an instance of the archive's gain data for $(D,L)$, $K$ a set of
gain constants and $C,\beta$ real. The package \emph{satisfies the lower link}
(\texttt{LinkG\_at}) if:
\begin{enumerate}
\item for every word $u$ and place $w$ whose residue characteristic is not
$T(u)$: a local length comparison (\texttt{LocLe}) at $w$ holds between the sublattice $\mathrm{NuLat}(u)$ recorded by $L$ and
the ramified lattice of $u$, with exponent
$\lfloor\operatorname{ur}_w(\operatorname{ta}_w)\rfloor$; the recorded length
at $(u,w)$ is at most $\dim(\text{ramified forms of }u)\cdot(\operatorname{ta}_w+\operatorname{io}_w)$;
and if the archive's set $\mathrm{VV}$ at the residue characteristic is empty,
then $\operatorname{ta}_w=\operatorname{io}_w=0$ and the recorded length is $0$;
\item for every prime $p$ and $t_0\in\mathrm{VV}(p)$, the archive's sum
$S_v(p,t_0)$ is at most $2m\,N_v(p,t_0)\,a/p$;
\item for every place $w$ whose residue characteristic $p$ is inert and
satisfies $p\le X''(n)$:
\begin{enumerate}
\item $0\le\operatorname{gain}_j(w)\le 4|TT|a/p$ for every lift index $j$;
\item with $TT_w\subseteq TT$ the set of top places attached to $w$ by $L$,
\[
\sum_{j=1}^m\sum_{t\in TT_w}\bigl(\mathrm{vM}(j,t,p)-\mathrm{lo}(p)\bigr)-m\,\mathrm{fl}(p)
\le\sum_{j=1}^m\operatorname{gain}_j(w)-\operatorname{ta}_w;
\]
\end{enumerate}
\item the archive's predicates \texttt{ServedOrbitsContained} and
\texttt{MeanKappaBound} hold.
\end{enumerate}
A family $\mathrm{Fm}$ satisfies \texttt{LinkG} if there are $C$, $x_0$ and a
function $\delta$ with $\delta(x)\to0$ such that every member $D(x,i)$ with
$x\ge x_0$ satisfies the lower link with the family's gain data and constants,
this $C$, and $\beta=b^{\delta(x)}$.
\end{definition}

\begin{remark}[what the lower link does and does not say]
Clause (3) is the finite-side input: at inert $p\le X''(n)$ the total gain is
bounded below by the Legendre-type sum, and each gain by $4|TT|a/p$. These are
hypotheses on the real numbers $\operatorname{gain}_j(w)$ of the size datum;
the archive's statement does not mention periods. That the gains are $p$-adic
valuations of lifted periods is the role assigned to them by the rest of the
package and was not traced here. The size datum's field \texttt{TTw\_full}
forces $TT_w=TT$ when $p$ has class $\mathrm{signal}$ or $p\ge X'$.
\end{remark}

\section{The upper link}

\begin{definition}[the Bessel identity]\label{geob:bessel}
\lean{Formalism.MainNotions.BesselId, Formalism.BasicNotions.D7.BuRational, Formalism.BasicNotions.D7.besS}
% uses-lean: Formalism.BasicNotions.D7, Formalism.BasicNotions.DivData
\uses{geoa:divdata} % resolved
The package $D$ satisfies the \emph{Bessel identity} if:
\begin{enumerate}
\item (rationality) for every word $u$, every $\varphi$ in the $\mathbb L$-space
of $N_u$ and every $y$ in that of $R$, the complex pairing
$B_u(\varphi,y)$ of the package's spaces equals the image under the fixed
embedding $\mathbb L\to\CC$ of the $\mathbb L$-valued pairing
$B_k(u)(\varphi)(y)$;
\item there is a unitary automorphism $V$ of the test space $\mathcal E$ such
that for every word $u$ there are $a_u\in\CC^{\times}$ and a unitary
automorphism $U_u$ of the source space $\mathcal M_u$ with
$B_u(\varphi,F)=a_u\operatorname{bes}_u(U_u\varphi,VF)$ for all $\varphi,F$,
where $\operatorname{bes}_u(\varphi,F)=\langle s(F),P_u(\varphi)\rangle$ is the
inner product of the spectral image of $F$ with the period image of $\varphi$.
\end{enumerate}
\end{definition}

\begin{definition}[vanishing above the cutoff]\label{geob:zeroabove}
\lean{Formalism.BasicNotions.SizeData₂.ZeroAbove}
\uses{geob:sz}
The size datum $L$ \emph{vanishes above the cutoff} if for every word $u$,
lift index $j$ and place $w$ whose residue characteristic $p$ satisfies
$p>X''(n)$, $p\ne T(u)$, $p\notin\{g_0,g_1\}$, $p\ne d$ and $p\nmid
\Gamma_u$ (the stabiliser bound of $u$), one has
$\operatorname{gain}_j(w)=E(u,j,w)=0$.
\end{definition}

\begin{definition}[the upper link]\label{geob:upperlink}
\lean{Formalism.MainNotions.LinkDIV_at}
\uses{geob:main, geob:sz, geob:bessel, geob:zeroabove}
% uses-lean: Formalism.BasicNotions.DivData.OneMapIntoCC, Formalism.BasicNotions.sFrame, Formalism.BasicNotions.D7.RRFull, Formalism.HN.htMap, Formalism.BasicNotions.SizeData₂.La, Formalism.BasicNotions.SizeData₂.B0, Formalism.BasicNotions.SizeData₂.Nc
\uses{geoa:divdata, geoa:slopes} % resolved
Let $R$ be the test lattice of $D$, $N_u$ the source lattices,
$B_k(u)\colon N_u\otimes\mathbb L\to R^{\vee}\otimes\mathbb L$ the period maps and
$h_u$ the recorded heights (the period lattice data of the preceding chapter).
Put $L_a=\log(a\operatorname{rd}_E)$ and let $B_0$, $N_c$ be the size budgets of
$L$. For real $C$, $C_{\mathrm{abs}}$, $C_H$, the size datum \emph{satisfies
the upper link} (\texttt{LinkDIV\_at}) if all of the following hold:
\begin{enumerate}
\item each $B_k(u)$ maps the lattice $N_u$ into the dual lattice of $R$;
\item $s_{\mathrm{frame}}+\mu_{\max}(R)\le x^{C_H}b$, where
$s_{\mathrm{frame}}=\sum_u\rk N_u\,(\mu(N_u)-h_u)/\sum_u\rk N_u$;
\item $L$ vanishes above the cutoff (Definition~\ref{geob:zeroabove});
\item the Bessel identity holds (Definition~\ref{geob:bessel});
\item $R$ is full: its lattice is finitely generated and spans $R\otimes\mathbb L$;
\item for every $u$ with $B_k(u)\ne0$, the height of $B_k(u)$ relative to the
hermitian structures of $N_u$ and $R^{\vee}$ is at most $h_u$;
\item for every word $u$,
\[
\mu(R^{\vee})+h_u-\mathrm{Fo}\le\mathcal M(u)+C\,|TT|\,a\,L_a+C\,m\,(B_0+N_c)\,x^{C_{\mathrm{abs}}\eps};
\]
\item for every word $u$ and lift index $j$,
$\bigl|\sum_{w\in S}\operatorname{sz}_w(E(u,j,w))-\bigl(\sum_{t}a_t\bigr)\log(g_0g_1)\bigr|\le CB_0$;
\item for all words $u,u'$,
$\bigl|\bigl(h_u+\sum_ta_t\log T(u)\bigr)-\bigl(h_{u'}+\sum_ta_t\log T(u')\bigr)\bigr|
\le C\bigl(|TT|aL_a+mB_0\bigr)$.
\end{enumerate}
\end{definition}

\begin{definition}[the upper link for a family]\label{geob:upperlink-fam}
\lean{Formalism.MainNotions.LinkDIV, Formalism.MainNotions.LinkDIVt}
\uses{geob:upperlink}
% uses-lean: Formalism.MainNotions.FamP, Formalism.MainNotions.LinkEta, Formalism.MainNotions.DIVinN
\uses{geoa:famcore} % resolved
A family $\mathrm{Fm}$ satisfies \texttt{LinkDIV} with exponents
$C_{\mathrm{abs}},C_H$ if there are $C,x_0$ such that every member with
$x\ge x_0$ satisfies the upper link with $C,C_{\mathrm{abs}},C_H$. The
predicate \texttt{LinkDIVt} is the implication: if the family satisfies
\texttt{LinkEta} and \texttt{DIVinN}, then it satisfies \texttt{LinkDIV}.
\end{definition}

\begin{remark}[the upper link and clause (2) of the record]
Clause (7) of the upper link is, word for word, the inequality
$\mu_R+h-\mathrm{Fo}\le\sum_t\bigl(m(\operatorname{ph}(a_t)-a_t\log\Omega)-a_t\log T\bigr)+\cdots$
that the slope record calls \texttt{DIV}
(\texttt{Formalism.RankGapBounds.WeightedMeanData.Data.DIV}); clause (8) has the
shape of the record's \texttt{EComp5} and clause (9) that of \texttt{HComp}.
The upper link is a hypothesis on the package; the archive's supply of it goes
through \texttt{LinkDIVt} and the covolume statements below, and is outside
this chapter.
\end{remark}

\section{Closed properties and the covolume}

\begin{definition}[thresholds and large families]\label{geob:threshold}
\lean{Formalism.MainNotions.TransportedAlt, Formalism.MainNotions.TransportedAlt.sup, Formalism.MainNotions.FamP.Large, Formalism.MainNotions.FamP.cPre, Formalism.MainNotions.FamP.cOut, Formalism.MainNotions.FamP.cMid, Formalism.MainNotions.FamP.cLast}
% uses-lean: Formalism.MainNotions.FamP
\uses{geoa:famcore} % resolved
A \emph{threshold} $K$ is a real number $K_{\mathrm{pre}}$ and three functions
$K_t,K_{\mathrm{mid}},K_{\mathrm{last}}$ from finite lists of reals to reals;
$K\vee K'$ is the pointwise maximum. A family $\mathrm{Fm}$ carries constant
lists $\mathrm{cPre}$ (one entry), $\mathrm{cOut}$ (sixteen entries),
$\mathrm{cLast}$ (thirty-two entries) and, for each real $c$, a list
$\mathrm{cMid}(c)$ of ten entries whose first entry is $c$. The family is
\emph{large for $(K,c)$} if every entry of $\mathrm{cPre}$ is at least
$K_{\mathrm{pre}}$, every entry of $\mathrm{cMid}(c)$ is at least
$K_{\mathrm{mid}}(\mathrm{cOut})$, and every entry of $\mathrm{cLast}$ is at
least $K_{\mathrm{last}}(\mathrm{cOut}\text{ followed by }\mathrm{cMid}(c))$.
\end{definition}

\begin{remark}[an unused component]
The second conjunct of \texttt{FamP.Large} is the literal \texttt{True}, so
the component $K_t$ of a threshold plays no role. In particular the entries
of $\mathrm{cOut}$ are not bounded below; they enter only as the argument of
$K_{\mathrm{mid}}$ and $K_{\mathrm{last}}$.
\end{remark}

\begin{definition}[closed properties]\label{geob:closed}
\lean{Formalism.MainNotions.FamClass, Formalism.MainNotions.Closed}
\uses{geob:threshold}
A class of families is a predicate $\mathcal C$ on families at every $\eps$.
Let $P(\mathrm{Fm},c)$ be a property of a family and a real number. Then $P$ is
\emph{closed on $\mathcal C$} if there is $\eps_1>0$ such that for every
$\eps\in(0,\eps_1]$ there is a threshold $K$ with: for every family
$\mathrm{Fm}$ at $\eps$ in $\mathcal C$ and every real $c$ such that
$\mathrm{Fm}$ is large for $(K,c)$, $P(\mathrm{Fm},c)$ holds.
\end{definition}

\begin{lemma}[closed properties combine]\label{geob:closed-and}
\lean{Formalism.MainNotions.FamP.large_sup, Formalism.MainNotions.Closed.and, Formalism.MainNotions.Closed.imp}
\uses{geob:closed}
A family is large for $(K\vee K',c)$ if and only if it is large for $(K,c)$
and for $(K',c)$. If $P$ and $Q$ are closed on $\mathcal C$, so is $P\wedge Q$.
If $P$ is closed on $\mathcal C$ and $P(\mathrm{Fm},c)$ implies
$Q(\mathrm{Fm},c)$ for every family at a positive $\eps$, then $Q$ is closed
on $\mathcal C$.
\end{lemma}
\begin{proof}
\uses{geob:threshold}
The first claim is the componentwise statement $\max(\alpha,\beta)\le\gamma$
iff $\alpha\le\gamma$ and $\beta\le\gamma$. For the conjunction take
$\eps_1=\min$ of the two values and, for each $\eps$, the threshold
$K\vee K'$; then use the first claim. The implication keeps $\eps_1$ and $K$.
\end{proof}

\begin{definition}[size of the field of values]\label{geob:fieldsmall}
\lean{GlobalTheory.FinalJoin2.FieldSmall, Formalism.MainNotions.WeightAlgebra.PO, Formalism.L2.Pj.FO}
% uses-lean: Formalism.BasicNotions.D7, Formalism.BasicNotions.SizeData₂
Let $F_O\subseteq\CC$ be the field generated over $\Q$ by the values of the
map \texttt{phiO} of the package, and put
$P_O=19\,f\,r^2L_x+198\,f^2L_x$ with $r=k+m-1$. The pair $(D,L)$ has a
\emph{small field of values} if $\log[F_O:\Q]\le P_O$ and
$\log|\operatorname{disc}F_O|/[F_O:\Q]\le P_O$.
\end{definition}

\begin{definition}[label tie of the $E$-series]\label{geob:esertie}
\lean{GlobalTheory.GLq.ESerTie}
% uses-lean: GlobalTheory.GLq.LabelDatum, GlobalTheory.GLq.eSerC
\uses{fam:green} % resolved
The package satisfies the \emph{label tie} if for $k=0,1$ there is a label
datum $\mathcal L_k$ for the auxiliary prime $g_k$ whose Teichmüller datum is the
one recorded by $D$ at $k$, and such that for every polynomial $f$ over
$\Z/g_k\Z$ the $E$-series of $D$ at $(k,f)$ equals the $E$-series of
$\mathcal L_k$ in rank $r$ at $f$.
\end{definition}

\begin{definition}[tied points]\label{geob:rectied}
\lean{GlobalTheory.CovTop.RecTied, GlobalTheory.CovTop.RowsT, GlobalTheory.CovTop.EvTied}
\uses{geob:fieldsmall, geob:esertie}
% uses-lean: GlobalTheory.BigTop.Pt, GlobalTheory.BigTop.PtAt
Put $\mathrm{RecTied}(D,L)$ for the conjunction of
Definitions~\ref{geob:fieldsmall} and~\ref{geob:esertie}. A point $p$ of a
member $D(x,i)$ of a family (the archive's \texttt{Pt}: a realisation of the
member by row data) is \emph{tied} (\texttt{RowsT}) if its rows hold and
$\mathrm{RecTied}(D(x,i),L(x,i))$. A property $Q$ of points holds
\emph{eventually at tied points} (\texttt{EvTied}) if there is $x_0$ such that
$Q$ holds at every tied point of every member $D(x,i)$ with $x>x_1$ and
$x\ge x_0$.
\end{definition}

\begin{remark}[\texttt{RecTied} is not the tie of a chain]
Despite its name, \texttt{RecTied} is a size bound on $F_O$ together with a
label condition on $E$-series. Both are conditions on the member's package and
size datum, not on the point $p$, and the predicate is unrelated to the tie
\texttt{ChainTied} of a family used in the last section.
\end{remark}

\begin{definition}[the covolume predicate]\label{geob:L3}
\lean{Formalism.L2.Indef.IndefData.L3, Formalism.L2.Indef.IndefScale.num, Formalism.L2.DIV.ZscOfNum, Formalism.L2.DIV.ZscOfNum3}
% uses-lean: Formalism.L2.Indef.IndefData, Formalism.L2.Indef.IndefNum, Formalism.L2.Indef.IndefScale
An indefinite datum $S$ has a degree $D_S$ and, for each lift index $j$, a
lattice $M_j$ in a space with archimedean Gram matrices $X_{W,j}$; let
$\operatorname{lcv}(M_j)=\frac2{[\mathbb L:\Q]}\log\operatorname{ht}(M_j)$ be its
normalised logarithmic covolume. For a numerical datum $N$ with bounds
$l_3(j)$, the predicate $\mathrm{L3}(S,N)$ is
$\operatorname{lcv}(M_j)/D_S\le l_3(j)$ for all $j$. A scale $Z$ gives, for
real $C,c$, the numerical datum $Z.\mathrm{num}(C,c)$ with
$l_3(j)=l_{3,\mathrm{main}}(j)+C\,l_{3,\mathrm{unit}}$ (and with $hh=x^cb$). For
the scale \texttt{ZscOfNum} built from a numerical datum, and its variant
\texttt{ZscOfNum3} which copies these two fields,
\[
l_{3,\mathrm{main}}(j)=\sum_{\rho}\bigl((k+j-1)-(m-j)\bigr)\,w(j,\rho)\,\log\Omega,
\qquad l_{3,\mathrm{unit}}=B_0\log\log x,
\]
the sum over the real places $\rho$ and $w(j,\rho)$ the weights.
\end{definition}

\begin{definition}[the covolume condition]\label{geob:covolx}
\lean{GlobalTheory.CovTop.CovolX}
\uses{geob:closed, geob:rectied, geob:L3}
% uses-lean: GlobalTheory.BigTop.FamHyp
The \emph{covolume condition} \texttt{CovolX} is the statement that the
following property of a family $\mathrm{Fm}$ and a real $c$ is closed on the
class \texttt{FamHyp}: there is $C_0\ge0$ such that, eventually at tied points
$p$ of $\mathrm{Fm}$, for every $C\ge C_0$ the covolume predicate
$\mathrm{L3}$ holds for the indefinite datum $p.S(\mathrm{false})$ and the
numerical datum $p.N(C,c)$.
\end{definition}

\begin{definition}[the covolume statement on the main chain]\label{geob:covolbound}
\lean{GlobalTheory.CovTop.RowsTAlt, GlobalTheory.CovTop.EventuallyAtAllPoints, GlobalTheory.CovTop.CovolumeBoundOfX}
\uses{geob:closed, geob:rectied, geob:L3}
% uses-lean: GlobalTheory.BigTop.FamHyp
\texttt{CovolumeBoundOfX} is \texttt{CovolX} with the condition ``the rows of
$p$ hold'' replaced by the archive's \texttt{PointConditions'} of $p$; the
condition $\mathrm{RecTied}$ is kept (\texttt{RowsTAlt}).
\end{definition}

\begin{theorem}[supply of the covolume statement]\label{geob:covolsupply}
\lean{SelfIntersectionVolumeEstimates.topCovolSlim2_use_alt_alt}
\uses{geob:covolbound}
% uses-lean: GlobalTheory.CompanionTheory.Hypotheses₂.ResolutionCharZero, GlobalTheory.CompanionTheory.Hypotheses₂.PrincipalisationCharZero, GlobalTheory.CompanionTheory.Hypotheses₂.PoincareLelongInv, PositivityTheory.MbHsUse.HSUse, ReciprocityTheory.ChartRestriction.G.ArakelovGeometry.MoriwakiBigStatement, GlobalTheory.Published.SemiAbelianLieSeminegative, GlobalTheory.CovTop.PublishedHodgeBundleTheorems
Assume the archive's statements of: resolution of singularities and
principalisation in characteristic zero; a Poincaré--Lelong formula; a use of
the arithmetic Hilbert--Samuel theorem (\texttt{HSUse}); Moriwaki's bigness
statement; an exactness statement of Fargues; Zarhin's statement; nef
Riemann--Roch; a Siegel toroidal integral; toroidal compactification for a fan;
seminegativity of the Lie algebra of a semi-abelian scheme; separation of
points; an arithmetic degree statement; the record of published theorems on
the Hodge bundle; and Grothendieck's deformation statement. Then
\texttt{CovolumeBoundOfX} holds.
\end{theorem}
\begin{proof}
% uses-lean: SelfIntersectionVolumeEstimates.topCovolG_alt_alt, PositivityTheory.HtPt.Ah.deep_slim_use, GlobalTheory.GreenBridge.green_of_regular
The proof is a single term: it assembles two hypothesis records from the
inputs (adding Green's parametrisation of cuspidal characters, proved in the
archive) and applies \texttt{topCovolG\_alt\_alt} together with
\texttt{deep\_slim\_use}. The development behind these lives under
\arch{Development/AutomorphicArithmetic/ThetaCorrespondence/ThetaPairLifts/HodgeBundleWeights/}
and was not expanded.
\end{proof}

\begin{remark}[where the covolume enters]
The lattices bounded by $\mathrm{L3}$ are the $M_j$ on the indefinite side, one
per lift index; they are not the test lattice $R$ of the upper link. The
coefficient $(k+j-1)-(m-j)$ is $p_j-q_j$ for the signature
$(p_j,q_j)=(k+j-1,m-j)$ of the $j$-th indefinite space, and the $\log\Omega$ in
$l_{3,\mathrm{main}}$ is the same period as in $\mathcal M(u)$. That the
covolume bounds summed over $j$ produce the term $-ma_t\log\Omega$ of the
upper link is an interpretation; the supply of the upper link was not traced.
\end{remark}

\section{From the slope record to the gap bound}

\begin{definition}[large scales and the error exponent]\label{geob:largex}
\lean{Formalism.LFunctionChainParameters.LErr.KBig, Formalism.LFunctionChainParameters.LErr.LargeX, Formalism.LFunctionChainParameters.LErr.CErr}
% uses-lean: Formalism.RankGapBounds.WeightedMeanData.Consts
\uses{geoa:record} % resolved
Put $K_{\mathrm{big}}(K,c_{HH},e_K)=200(1+K)^2(1+c_{HH})(1+e_K)$. A real $x$ is
\emph{large} for $(\eps,K,c_{HH},e_K)$ if $\log x\le x^{\eps}$ and
$K_{\mathrm{big}}(K,c_{HH},e_K)\le x^{\eps}$. For record constants $C$ put
$C_{\mathrm{Err}}=e_{\mathrm{Div}}+e_{\mathrm{Loss}}+3$.
\end{definition}

\begin{definition}[the pointwise error and the pointwise bound]\label{geob:pointwise}
% overlaps (deduplicated): geoa:clauses-5
\lean{Formalism.LFunctionChainParameters.WeightedSlopeBounds.MLo5, Formalism.LFunctionChainParameters.WeightedSlopeBounds.UBar5, Formalism.LFunctionChainParameters.WeightedSlopeBounds.ErrX5at, Formalism.LFunctionChainParameters.WeightedSlopeBounds.PointwiseBoundVariant}
% uses-lean: Formalism.RankGapBounds.WeightedMeanData.Data, Formalism.RankGapBounds.WeightedMeanData.Data.HidaSetting, Formalism.RankGapBounds.WeightedMeanData.Data.DIV, Formalism.RankGapBounds.WeightedMeanData.Data.HComp, Formalism.RankGapBounds.WeightedMeanData.Data.Slopes, Formalism.RankGapBounds.WeightedMeanData.Data.SlopeIneq, Formalism.RankGapBounds.WeightedMeanData.Data.HHBound, Formalism.RankGapBounds.WeightedMeanData.Data.FParams, Formalism.RankGapBounds.WeightedMeanData.Data.Vvar, Formalism.RankGapBounds.WeightedMeanData.Data.LoP
\uses{geoa:clauses, geoa:record, geoa:record-derived} % resolved
Let $\mathcal D$ be a slope record with index set $U_0$, word sets $W(u_0)$,
constants $C$, and a real $X''$. With the record's fields, put
\[
\mathrm{Lo}_5(u_0,u)=c_{Sl1}\mathrm{IO}(u_0)+c_{Sl1}n_Tm\Bigl(\frac ak+b'\Bigr)\log T(u_0,u)+c_{Sl1}X''+1 .
\]
The variants of the record clauses used here are: \texttt{QuotientSlopeBoundsWith}
($\mu_N^{\max}\le\mu_M^{\max}+\mathrm{Ga}+\mathrm{Fo}-\operatorname{Ee}+c_{Sl1}mB_0$ and
$\mu_M+\mathrm{Ga}+\mathrm{Fo}-\operatorname{Ee}-\mathrm{Lo}_5\le\mu_N$ for all
$u_0,u$); \texttt{EComp5}
($|\operatorname{Ee}(u_0,u)-m\,(\sum_ta_t)\lg|\le c_Em B_0$); \texttt{Count5}
(the lower bound
$\sum_j\frac1f\sum_{\rho}\log(w(u_0,j,\rho)v/\pi)+m\lg-c_{\mathrm{Count}}(L_a+m^2L_a/k)\le\log T(u_0,u)$,
the upper bound $\log T\le c_{\mathrm{Count}}mL_x$, and
$|\log T(u_0,u)-\log T(u_0,u')|\le c_{\mathrm{Count}}\log k$); and
\texttt{LossBounds} ($\mathrm{PS}(u_0)-\mathrm{Ga}(u_0)\le\mathrm{Lo}_P$ and
$\mathrm{IO}(u_0)\le\mathrm{Lo}_P$ for every $u_0$). The quantity
$\mathrm{Err}_5(u_0)$ (\texttt{ErrX5at}) is the explicit sum
\begin{multline*}
\textstyle(\sum_ta_t)c_{\mathrm{Count}}\bigl(L_a+\frac{m^2L_a}k\bigr)+c_{\mathrm{Div}}n_TaL_a+c_{\mathrm{Div}}m(B_0+N_c)x^{e_{\mathrm{Div}}\eps}
+mc_EB_0+c_{\mathrm{Slope}}n_TaL_a\bigl(1+\frac mk\bigr)\\
\textstyle+\log|W(u_0)|+\mathrm{Lo}_P+\overline{\mathrm{Lo}}_5
+(1-\bar\delta)^{-1}\bigl(4\log(1+x^{c_{HH}})(\bar U_5+b)+(\bar\eta+4\sqrt{\bar e})x^{c_{HH}}b\bigr),
\end{multline*}
where $\overline{\mathrm{Lo}}_5$ is $\mathrm{Lo}_5$ with $\mathrm{IO}$ replaced by
$\mathrm{Lo}_P$ and $\log T$ by $c_{\mathrm{Count}}mL_x$, and
$\bar U_5=c_{\mathrm{Slope}}n_TaL_a(2+m/k)+V_{\mathrm{var}}+c_{Sl1}mB_0+\overline{\mathrm{Lo}}_5$.
The \emph{pointwise bound} for $(\mathcal D,C,X'')$ is the statement: if the
record satisfies \texttt{HidaSetting}, \texttt{DIV}, \texttt{EComp5},
\texttt{HComp}, \texttt{Count5}, \texttt{Slopes},
\texttt{QuotientSlopeBoundsWith}, \texttt{LossBounds}, \texttt{SlopeIneq},
\texttt{HHBound} and \texttt{FParams}, then
$\mathrm{PS}(u_0)\le\operatorname{Main}(u_0)+\mathrm{Err}_5(u_0)$ for every
$u_0$.
\end{definition}

\begin{lemma}[lower bound for the mean source slope]\label{geob:slower}
\lean{Formalism.LFunctionChainParameters.WeightedSlopeBounds.s_lower5, Formalism.LFunctionChainParameters.WeightedSlopeBounds.K5v5, Formalism.LFunctionChainParameters.WeightedSlopeBounds.Lam5, Formalism.LFunctionChainParameters.WeightedSlopeBounds.exponent_le_upper}
\uses{geob:pointwise, geoa:clauses, geoa:clauses-5}
% uses-lean: Formalism.RankGapBounds.WeightedMeanData.Data.s, Formalism.RankGapBounds.WeightedMeanData.wmean
\uses{geoa:record-derived} % resolved
Assume the record satisfies \texttt{HidaSetting}, \texttt{DIV}, \texttt{EComp5},
\texttt{Count5}, \texttt{Slopes} and \texttt{QuotientSlopeBoundsWith}. Then for
every $u_0$,
\[
\mathrm{Ga}(u_0)+\mu_R(u_0)-K_5(u_0)-\operatorname{wmean}_{d(u_0,\cdot)}\mathrm{Lo}_5(u_0,\cdot)\le s(u_0),
\]
where $s(u_0)$ is the $d$-weighted mean of $\mu_N(u_0,u)-h(u_0,u)$ over $u$ and
$K_5(u_0)=c_{\mathrm{Slope}}n_TaL_a(1+m/k)+m\bigl((\sum_ta_t)\lg+c_EB_0\bigr)+\mathrm{SDiv}-(\sum_ta_t)\Lambda_5(u_0)+\mathrm{EDiv}$,
with $\Lambda_5(u_0)$ the lower bound for $\log T$ in \texttt{Count5},
$\mathrm{SDiv}=\sum_t m(\operatorname{ph}(a_t)-a_t\log\Omega)$ and
$\mathrm{EDiv}$ the error of \texttt{DIV}.
\end{lemma}
\begin{proof}
It suffices to bound $\mu_N(u_0,u)-h(u_0,u)$ from below for each word $u$ and
average. The second inequality of \texttt{QuotientSlopeBoundsWith} bounds
$\mu_N$ below by $\mu_M+\mathrm{Ga}+\mathrm{Fo}-\operatorname{Ee}-\mathrm{Lo}_5$;
\texttt{DIV} bounds $-h$ below by $\mu_R-\mathrm{Fo}$ minus the main term and its
error; the main term contains $-(\sum_ta_t)\log T$, bounded using the lower
bound of \texttt{Count5} (with $\sum_ta_t\ge0$ from \texttt{HidaSetting});
\texttt{EComp5} bounds $\operatorname{Ee}$ above; \texttt{Slopes} bounds $\mu_M$
below. The terms $\mathrm{Fo}$ cancel and the rest is linear arithmetic.
\end{proof}

\begin{lemma}[the gain is at most the main term]\label{geob:gabound}
\lean{Formalism.LFunctionChainParameters.WeightedSlopeBounds.Ga_bound5, Formalism.LFunctionChainParameters.WeightedSlopeBounds.K0v5, Formalism.LFunctionChainParameters.WeightedSlopeBounds.key_identity5, Formalism.LFunctionChainParameters.WeightedSlopeBounds.u_function_le_mean, Formalism.LFunctionChainParameters.WeightedSlopeBounds.Vvar_bound5}
\uses{geob:slower, geob:pointwise}
% uses-lean: Formalism.RankGapBounds.WeightedMeanData.Data.key_identity, Formalism.RankGapBounds.WeightedMeanData.Data.uF, Formalism.RankGapBounds.WeightedMeanData.Data.K3
\uses{geoa:record-derived} % resolved
Assume in addition \texttt{HComp}, \texttt{SlopeIneq}, \texttt{HHBound} and
\texttt{FParams}. Then for every $u_0$,
\[
\mathrm{Ga}(u_0)\le\operatorname{Main}(u_0)+\log|W(u_0)|+K_0+\overline{w}
+(1-\bar\delta)^{-1}\Bigl(4\log(1+x^{c_{HH}})\bigl(K_3+\overline{w}+b\bigr)+(\bar\eta+4\sqrt{\bar e})x^{c_{HH}}b\Bigr),
\]
with $\overline{w}=\operatorname{wmean}_{d(u_0,\cdot)}\mathrm{Lo}_5(u_0,\cdot)$,
$K_0=(\sum_ta_t)c_{\mathrm{Count}}(L_a+m^2L_a/k)+\mathrm{EDiv}+mc_EB_0+c_{\mathrm{Slope}}n_TaL_a(1+m/k)$
and $K_3$ the record's constant.
\end{lemma}
\begin{proof}
\texttt{SlopeIneq}, applied with the bounds of \texttt{FParams} and
\texttt{HHBound}, gives
$s-\log|W|-(1-\delta)^{-1}\bigl(4u_F\log(1+HH/u_F)+(\eta+4\sqrt e)HH\bigr)\le\mu_R$,
where $u_F=\max_u(\mu_N^{\max}-h)-s\ge0$. Lemma~\ref{geob:slower} gives
$s\ge\mathrm{Ga}+\mu_R-K_5-\overline{w}$, so $\mu_R$ cancels. The archive bounds
$u_F\le K_3+\overline w$ (\texttt{u\_function\_le\_mean}, from both halves of
\texttt{QuotientSlopeBoundsWith}, \texttt{Slopes}, and the variation bound
\texttt{Vvar\_bound5} built on \texttt{EComp5}, \texttt{HComp} and
\texttt{Count5}), then bounds $4u_F\log(1+HH/u_F)$ through \texttt{phi\_bound}
using $HH\le x^{c_{HH}}b$, replaces $\delta,\eta,e$ by their bars, and uses the
identity \texttt{key\_identity5}, which rewrites
$\mathrm{SDiv}-(\sum_ta_t)\Lambda_5+m(\sum_ta_t)\lg$ as
$\operatorname{Main}(u_0)+(\sum_ta_t)c_{\mathrm{Count}}(L_a+m^2L_a/k)$. The
remaining step is linear arithmetic.
\end{proof}

\begin{remark}[the test lattice cancels]
The slope $\mu_R$ of $R^{\vee}$ enters Lemma~\ref{geob:slower} with sign $+$
(through the upper bound \texttt{DIV}) and \texttt{SlopeIneq} with sign $-$, so
it does not appear in Lemma~\ref{geob:gabound}. Neither do the heights $h$. The
two lattice inequalities therefore act only as a sandwich for the one number
$\mu_R$.
\end{remark}

\begin{theorem}[the pointwise bound]\label{geob:eachunit}
\lean{Formalism.LFunctionChainParameters.WeightedSlopeBounds.each_unit_bound_statement, Formalism.LFunctionChainParameters.WeightedSlopeBounds.weighted_mean_lower_bound_le}
\uses{geob:pointwise, geob:gabound}
For every slope record $\mathcal D$, constants $C$ and real $X''$ the pointwise
bound of Definition~\ref{geob:pointwise} holds: under the eleven listed
clauses, $\mathrm{PS}(u_0)\le\operatorname{Main}(u_0)+\mathrm{Err}_5(u_0)$ for
every $u_0$.
\end{theorem}
\begin{proof}
\texttt{LossBounds} gives $\mathrm{PS}\le\mathrm{Ga}+\mathrm{Lo}_P$ and
$\mathrm{IO}\le\mathrm{Lo}_P$. The weighted mean of $\mathrm{Lo}_5$ is at most
$c_{Sl1}\mathrm{IO}+c_{Sl1}n_Tm(a/k+b')c_{\mathrm{Count}}mL_x+c_{Sl1}X''+1$ by the
upper bound for $\log T$ in \texttt{Count5}, hence at most
$\overline{\mathrm{Lo}}_5$. Insert this into Lemma~\ref{geob:gabound}; the
coefficient $(1-\bar\delta)^{-1}4\log(1+x^{c_{HH}})$ is nonnegative since
$\bar\delta<1$. Linear arithmetic gives the claim.
\end{proof}

\begin{lemma}[size of the pointwise error]\label{geob:errx5}
\lean{Formalism.LFunctionChainParameters.LChain.errX5at_le, Formalism.LFunctionChainParameters.LErr.ErrText}
\uses{geob:pointwise, geob:largex}
% uses-lean: Formalism.LFunctionChainParameters.LErr.ErrWindows
\uses{geoa:errwindows} % resolved
Assume the record satisfies \texttt{ErrWindows} with $(C,X'',e_K,K)$ and that
$x$ is large for $(\eps,K,c_{HH},e_K)$. Then for every $u_0$,
\[
\mathrm{Err}_5(u_0)\le x^{C_{\mathrm{Err}}\eps}\Bigl(n_Ta\bigl(1+\tfrac{m^2}k\bigr)+m(B_0+N_c)\Bigr).
\]
\end{lemma}
\begin{proof}
Not expanded. The proof
(\arch{Formalism/SizeWindowBrackets/ChainSlopeWindows/PChain.lean}, from line
613) bounds each summand of $\mathrm{Err}_5$ using the size conditions in
\texttt{ErrWindows} and the two inequalities of largeness.
\end{proof}

\begin{theorem}[the gap bound]\label{geob:chtop}
% overlaps (deduplicated): geoa:main-eq
\lean{Deduction.ch_top}
\uses{geob:eachunit, geob:errx5, geob:largex}
% uses-lean: Deduction.CH.IdentU, Deduction.CH.BNBoundsU, Deduction.CH.gapA, Formalism.LFunctionChainParameters.LChain.SizeData₂, Formalism.LFunctionChainParameters.LErr.ErrWindows, Formalism.LFunctionChainParameters.LErr.errorBracketSizeWindows, Formalism.LFunctionChainParameters.characterSumData.MAIN
\uses{bal:gap, bal:main-defect, bal:rest-size, geoa:errwindows, geoa:sizedata} % resolved
Let $H$ be a finite abelian group, $\mathcal D$ a slope record with index set
$U_0$ (nonempty word sets $W(u_0)$), $\operatorname{tr}\colon U_0\to H$ a map,
$C$ constants, $X'',e_K,K$ real, and let $(D,\ell,N,TT,T_2,a,w)$ be arithmetic
data on $H$ with bijections $e_T$ from the top index set onto $TT$ and $e_P$
from $\{1,\dots,f\}$ onto the set $\Phi$ of places. Assume
\begin{enumerate}
\item the identification \texttt{IdentU} along $\operatorname{tr}$ (the record's
weights, $\gamma$, $L'/L$ value and $\mathrm{PS}$ are those of the arithmetic
data);
\item the eleven record clauses (\texttt{LChain.SizeData₂});
\item \texttt{ErrWindows} with $(C,X'',e_K,K)$;
\item \texttt{BNBoundsU} for $T_2$:
$B_0\le2|H|mbx^{\eps}$ and $N_c\le|T_2|kbx^{\eps}/(|H|^2x^{-\eps})$;
\item $x$ is large for $(\eps,K,c_{HH},e_K)$.
\end{enumerate}
Then for every $u\in U_0$,
\[
\operatorname{gap}(\operatorname{tr}u)\le x^{(e_{\mathrm{Div}}+e_{\mathrm{Loss}}+3)\eps}\,
\mathcal E(x,\eps,|H|,|TT|,|T_2|,m,b,k,a),
\]
where $\operatorname{gap}(h)=\sum_{t\in TT}\sum_j\bigl(\mathrm{ASK}-\Re\,\mathrm{MAIN}\bigr)(a_t,w(h,j))$
and $\mathcal E=n_Ta(1+m^2/k)+m\bigl(2nmbx^{\eps}+|T_2|kbx^{\eps}/(n^2x^{-\eps})\bigr)$
with $n=|H|$, $n_T=|TT|$.
\end{theorem}
\begin{proof}
By \texttt{ch\_top\_main\_eq}, \texttt{IdentU} gives
$\operatorname{Main}(u)=\sum_{t\in TT}\sum_j\Re\,\mathrm{MAIN}(\ell,a_t,w(\operatorname{tr}u,j))$,
and the last clause of \texttt{IdentU} gives $\mathrm{PS}(u)=\sum_t\sum_j\mathrm{ASK}$;
hence $\operatorname{gap}(\operatorname{tr}u)=\mathrm{PS}(u)-\operatorname{Main}(u)$.
Theorem~\ref{geob:eachunit} bounds this by $\mathrm{Err}_5(u)$, and
Lemma~\ref{geob:errx5} by $x^{C_{\mathrm{Err}}\eps}(n_Ta(1+m^2/k)+m(B_0+N_c))$.
Finally \texttt{BNBoundsU} and $n_T=|TT|$ (from $e_T$) give
$n_Ta(1+m^2/k)+m(B_0+N_c)\le\mathcal E$, monotonically.
\end{proof}

\begin{remark}[the size of the bracket]
Divide $\mathcal E$ by $n_Tmb$. Under \texttt{RestSize}
($a\le2xb$, $n_T\ge x^{-\eps/8}n/2$, $m\ge x^{1/2-\eps}$, $k\le n^2x^{1/2-\eps}$)
and the bounds $xm\le k$, $m\le x^{1/2+\eps/2}$ of \texttt{ErrWindows}, one
gets $\mathcal E/(n_Tmb)\ll(x/m+m+x^{1/2})x^{O(\eps)}$, whose largest term is
smallest at $m\approx x^{1/2}$. This is where the exponent $1/2$ of the
Chebyshev inequality comes from. The constraints on $m$ and $k$ are hypotheses
of the record, not consequences.
\end{remark}

\section{The Chebyshev inequality on windows}

\begin{definition}[core constants at a host]\label{geob:coreath}
% overlaps (deduplicated): bal:to-E
\lean{UpperTheory.Comparison₂.CoreAtH, UpperTheory.Comparison₂.FamCoreConstantBounds}
% uses-lean: Formalism.HC.HostRaw, Formalism.MainNotions.FamCore, Deduction.CH.IdentU, Deduction.CH.BNBoundsU, Deduction.CH.SizesA, Formalism.LFunctionChainParameters.LChain.SizeData₂, Formalism.LFunctionChainParameters.LErr.ErrWindows, Formalism.LFunctionChainParameters.LChain.CostWindows, Formalism.LFunctionChainParameters.LChain.RestSize, Formalism.LFunctionChainParameters.ThickSetMeanBounds.ThinScales, UpperTheory.Comparison₂.arith
\uses{bal:rest-size, bal:scales, bal:sizes, bal:to-E, geoa:errwindows, geoa:famcore, geoa:sizedata, host:host-raw} % resolved
Let $H$ be a raw host on a frame $\varphi$ with group $A$, and $TT\subseteq A$.
The host has \emph{core constants} $(e_D,e_L,K,c_{HH},e_K)$ at $(\eps,x)$
(\texttt{CoreAtH}) if there are a slope record $\mathcal D$ with index set $U_0$
and a map $\operatorname{tr}\colon U_0\to A$ whose image contains the subgroup
$A_0$, constants $C$, and data $\ell,N,T_2,a,w,e_T,e_P,\eta,X',\vartheta$ such
that:
\begin{enumerate}
\item $\ell(\psi)=L'(1,\chi)/L(1,\chi)$ for every nontrivial additive character
$\psi$ of $\{\pm1\}\times A$ and every primitive Dirichlet character $\chi$
whose values at $n\ne0$ are the values $\operatorname{pv}(\psi,n)$ given by the
host's Artin map;
\item the record has $x$ and $\eps$ as its scale, $e_{\mathrm{Div}}\le e_D$,
$e_{\mathrm{Loss}}\le e_L$ and $c_{HH}^{\mathcal D}\le c_{HH}$;
\item \texttt{IdentU} along $\operatorname{tr}$ for the arithmetic data of the
host, the eleven record clauses, \texttt{ErrWindows} with $(C,X'',e_K,K)$,
\texttt{BNBoundsU} for $T_2$;
\item the size conditions \texttt{SizesA}, the thin-scale bound
\texttt{ThinScales}, \texttt{CostWindows} and \texttt{RestSize};
\item $T_2\subseteq A_0$.
\end{enumerate}
A family core has \emph{core constants} at $x$ if some member $D(x,i)$ has a host
with core constants at $(\eps,x)$ for its top set $TT$
(\texttt{FamCoreConstantBounds}).
\end{definition}

\begin{lemma}[a tie gives core constants]\label{geob:tiedcore}
\lean{UpperTheory.Comparison₂.tiedCore_of_chainTied, UpperTheory.Comparison₂.arithOf_eq}
\uses{geob:coreath}
% uses-lean: Formalism.MainNotions.ChainTied, Formalism.MainNotions.FamP
\uses{geoa:famcore, geoa:tied} % resolved
Let $\mathrm{Fm}$ be a family at $\eps$ and suppose it is tied at $x$ with
constants $(e_D,e_L,K,c_{HH},e_K,C_L)$ (\texttt{ChainTied}). Then its core has
core constants $(e_D,e_L,K,c_{HH},e_K)$ at $x$.
\end{lemma}
\begin{proof}
Take the index $i_0$ of the tie, $U_0=A$ with $\operatorname{tr}$ the identity,
and the tie's record and data. The tie's $TT$ equals the host's top set by the
first arithmetic tie. The tie's clauses \texttt{Ident}, \texttt{BNBounds} and
the record clauses are the $\operatorname{tr}=\mathrm{id}$ cases of
\texttt{IdentU}, \texttt{BNBoundsU} and the same clauses; the size and
thin-scale parts of \texttt{GrowthInequalityRemainingHyps} give \texttt{SizesA}
and \texttt{ThinScales}; and $T_2\subseteq A_0$ follows from the second
arithmetic tie, since each element of $T_2$ is a digit, which lies in the cyclic
subgroup. The point, global and local ties, the weight tie, the clause
\texttt{ValuesAtX} and the constant $C_L$ are discarded.
\end{proof}

\begin{remark}[the Chebyshev inequality forgets the lattices]
Lemma~\ref{geob:tiedcore} drops \texttt{TiePts}, \texttt{TieGlobal} and
\texttt{TieLocal}, which identify the record's ranks, slopes, heights and local
sums with invariants of actual lattices and size data. Everything after this
point uses only inequalities among real numbers and the arithmetic identification
of $\ell$ with $L'/L(1,\chi)$. The lattice content of a tie matters only for
whether tied families exist, which is the statement \texttt{TopG} below.
\end{remark}

\begin{theorem}[the Chebyshev inequality at a host]\label{geob:coregm}
\lean{UpperTheory.Comparison₂.core_gm}
\uses{geob:coreath, geob:chtop, geob:largex}
% uses-lean: Deduction.ch_core, UpperTheory.Comparison₂.gmA_eq, UpperTheory.Comparison₂.sigA_le, UpperTheory.Comparison₂.charsA, UpperTheory.Comparison₂.hostFactsS, Formalism.LFunctionChainParameters.LErr.large_x_eventually, Formalism.HC.EAtH, Formalism.HC.Level.Gm
\uses{an:level, bal:ch-core, bal:charsA, bal:gmA-eq, bal:host-split, bal:sigA-le} % resolved
Let $\Lambda$ be a level and $e_D,e_L,K,c_{HH},e_K,C$ real, $\eps>0$. There is
$x_1$ such that for every $x\ge x_1$, every frame $\varphi$, every host $H$ on
$\varphi$ and every $TT\subseteq A$: if $H$ satisfies the explicit-formula
condition \texttt{EAtH} for $\Lambda$ at $(x,\eps,C)$ and has core constants
$(e_D,e_L,K,c_{HH},e_K)$ at $(\eps,x)$, then
\[
G_{\Lambda}(x)\le x^{1/2+(e_D+e_L+7)\eps}.
\]
\end{theorem}
\begin{proof}
Choose $x_1$ beyond three thresholds: that of \texttt{ch\_core} with exponent
$e_D+e_L+3$, that of largeness of $x$ for $(\eps,K,c_{HH},e_K)$, and one past
which $x^{\eps}\ge1+4o_{\Lambda}$. Given the core data, \texttt{ch\_core} is applied
with the character facts of the host (\texttt{charsA}, from \texttt{EAtH}) and
with the gap bound of Theorem~\ref{geob:chtop} at each element of $A_0$ (each is
$\operatorname{tr}u$ for some $u$), where
$x^{C_{\mathrm{Err}}\eps}\le x^{(e_D+e_L+3)\eps}$ since
$e_{\mathrm{Div}}\le e_D$ and $e_{\mathrm{Loss}}\le e_L$. After rewriting with
\texttt{gmA\_eq} and \texttt{sigA\_le} this gives
$G_{\Lambda}(x)\le x^{1/2+(e_D+e_L+6)\eps}+4o_{\Lambda} x^{1/2+\eps}$ up to the form
of the constants, and $x^{\eps}\ge1+4o_{\Lambda}$ absorbs the second term into the
extra factor $x^{\eps}$.
\end{proof}

\begin{definition}[Chebyshev inequality on windows]\label{geob:chebdefs}
\lean{Formalism.HC.ChebyshevInequalityOnFamilyWindows, Formalism.HC.ChebyshevInequalityOnWindows, Formalism.HC.ChebyshevInequalityAtWindowScale}
% uses-lean: Formalism.HC.EFamC, Formalism.MainNotions.FamCore, Formalism.MainNotions.FamCore.Windows, Formalism.HC.Level.Gm
\uses{an:level, geoa:efamc, geoa:famcore} % resolved
Fix a type of families $\mathrm{Fam}_{\eps}$ with a core map to family cores and a
tie predicate $\mathrm{Tied}(\mathrm{Fm},e_D,e_L,K,c_{HH},e_K,C_L,x)$.
\begin{enumerate}
\item $\mathrm{Fm}$ is \emph{tied on its windows} with $(e_D,e_L)$ if there are
$K,c_{HH},e_K,C_L,x_0$ such that $\mathrm{Tied}$ holds at every $x\ge x_0$ in
the window set $\mathrm{Win}$ of its core.
\item The \emph{Chebyshev inequality on windows} for a level $\Lambda$ is: for
all real $e_D,e_L$ and $\eps>0$, and every family $\mathrm{Fm}$ at $\eps$ whose
core satisfies \texttt{EFamC} for $\Lambda$ and which is tied on its windows
with $(e_D,e_L)$, there is $x_1$ with
$G_{\Lambda}(x)\le x^{1/2+(e_D+e_L+7)\eps}$ for every $x\ge x_1$ in $\mathrm{Win}$.
\item For $\lambda$ real, \texttt{ChebyshevInequalityAtWindowScale}$(\lambda,\Lambda)$
is: there are $e_D,e_L$ and $\eps_1>0$ such that for every $\eps\in(0,\eps_1)$
there is a family at $\eps$ whose core satisfies \texttt{EFamC} for $\Lambda$,
has windows of ratio $\lambda$, and is tied on its windows with $(e_D,e_L)$.
\end{enumerate}
\end{definition}

\begin{remark}[a misleading name]
Despite its name, \texttt{ChebyshevInequalityAtWindowScale} contains no
inequality for $G_{\Lambda}$: besides the order conditions $0<\eps_1$, $\eps<\eps_1$
and $x\ge x_0$, it asserts the existence of a family satisfying \texttt{EFamC}, with
windows, that is tied. On the main chain the tie is \texttt{ChainTied}, which does contain
inequalities between real numbers (for instance $e_{\mathrm{Div}}\le e_D$).
The inequality is \texttt{ChebyshevInequalityOnWindows}. The archive's Landau
step (\texttt{level\_of\_tiedWin}) combines the two, choosing $\eps$ so small
that $(e_D+e_L+7)\eps$ is below half the distance of $\Re s$ from $1/2$.
\end{remark}

\begin{lemma}[core constants give the Chebyshev inequality]\label{geob:windowineq}
\lean{UpperTheory.Comparison₂.windowInequality_of_artinInputs}
\uses{geob:coreath, geob:coregm, geob:chebdefs}
Let the tie predicate be such that a tie at $x$ with
$(e_D,e_L,K,c_{HH},e_K,C_L)$ gives core constants
$(e_D,e_L,K,c_{HH},e_K)$ at $x$ for the core. Then the Chebyshev inequality on
windows holds for every level $\Lambda$.
\end{lemma}
\begin{proof}
Given $\mathrm{Fm}$ with \texttt{EFamC} constants $(C,x_E)$ and tie constants
$(K,c_{HH},e_K,C_L,x_0)$, take $x_1$ from Theorem~\ref{geob:coregm} for
$(e_D,e_L,K,c_{HH},e_K,C,\eps)$ and put $x_1'=\max(x_1,x_0,x_E)$. For
$x\ge x_1'$ in $\mathrm{Win}$ the tie gives a member with core constants, and
\texttt{EFamC} gives the explicit-formula condition for that member's host.
Apply Theorem~\ref{geob:coregm}.
\end{proof}

\begin{theorem}[the Chebyshev inequality for tied families]\label{geob:levelgm}
\lean{UpperTheory.Comparison₂.level_gm_le_on_family_scales}
\uses{geob:windowineq, geob:tiedcore}
% uses-lean: Formalism.MainNotions.FamP, Formalism.MainNotions.ChainTied
\uses{geoa:famcore, geoa:tied} % resolved
For every level $\Lambda$, the Chebyshev inequality on windows holds for the
families \texttt{FamP} with their cores and the tie \texttt{ChainTied}.
Explicitly: for all $e_D,e_L$, $\eps>0$ and every family $\mathrm{Fm}$ at $\eps$
satisfying \texttt{EFamC} for $\Lambda$ and \texttt{ChainTied} with $(e_D,e_L)$
at every large $x$ in $\mathrm{Win}$, there is $x_1$ with
$G_{\Lambda}(x)\le x^{1/2+(e_D+e_L+7)\eps}$ for all $x\ge x_1$ in $\mathrm{Win}$.
\end{theorem}
\begin{proof}
Lemma~\ref{geob:windowineq} applied with Lemma~\ref{geob:tiedcore}.
\end{proof}

\begin{definition}[the existence statement \texttt{TopG}]\label{geob:topg}
\lean{UpperTheory.Top.FamilyClassConnected, UpperTheory.Top.TopG}
% uses-lean: UpperTheory.Top.FReq, UpperTheory.Top.FReq.cls, UpperTheory.Top.FReq.req, UpperTheory.Top.HypV, GlobalTheory.EX.HostExW, Formalism.MainNotions.FamClass, Formalism.MainNotions.ChainTied, Formalism.MainNotions.FamP
\uses{geoa:famcore, geoa:tied, geob:closed, host:gchi, host:host-exw} % resolved
A class $\mathcal E$ of families is \emph{connected}
(\texttt{FamilyClassConnected}) if there are $e_D,e_L$ and $\eps_1>0$ such that
for every $\eps\in(0,\eps_1)$ there are a family $\mathrm{Fm}$ at $\eps$ in
$\mathcal E$ and $K,c_{HH},e_K,C_L,x_0$ such that $\mathrm{Fm}$ is tied
(\texttt{ChainTied}) with these constants at every $x\ge x_0$ in its window set.
\texttt{TopG} is the statement: for every $\lambda>1$ and every host requirement
$G$, if hosts meeting $G$ exist on a system of windows of ratio $\lambda$
(\texttt{HostExW}$(\lambda,G)$), and every family in the class
$G.\mathrm{cls}(\lambda)$ satisfies the Sha bound \texttt{HypV}, then the class
$G.\mathrm{cls}(\lambda)$ is connected. Here $G.\mathrm{cls}(\lambda)$ consists
of the families with windows of ratio $\lambda$ all of whose members beyond the
family threshold meet $G$.
\end{definition}



\begin{remark}[what remains]
After Theorem~\ref{geob:levelgm}, the whole content of the geometric
inequality is the hypothesis \texttt{TopG}: the existence, at every small
$\eps$, of a family on windows that is tied at every large $x$ in its windows.
The tie requires the slope record clauses, in particular \texttt{DIV}
(the upper link), \texttt{SlopeIneq} and the losses, to hold for the actual
lattices and size data of the members. The archive's proof of \texttt{TopG}
goes through the record of hypotheses \texttt{Hyps}, in which the links of this
chapter and the covolume condition are inputs; it is outside this chapter.
\end{remark}
